\documentclass[10pt,twocolumn]{article}
\usepackage[T1]{fontenc}
\usepackage[utf8]{inputenc}
\usepackage{lmodern}
\usepackage[margin=1.8cm,top=2.4cm,bottom=2cm,columnsep=0.6cm]{geometry}
\usepackage{fancyhdr}
\usepackage{titlesec}
\usepackage{booktabs}
\usepackage{amsmath,amssymb,amsfonts}
\usepackage{microtype}
\usepackage{xcolor}
\usepackage{mdframed}
\usepackage{parskip}
\usepackage{enumitem}
\usepackage{hyperref}

\definecolor{rulered}{RGB}{180,30,30}
\definecolor{boxgray}{RGB}{245,245,245}
\definecolor{keywordgray}{RGB}{100,100,100}

\pagestyle{fancy}
\fancyhf{}
\fancyhead[L]{\textsc{\small Claude Mag}}
\fancyhead[C]{\small\textit{ICML 2026 Special Edition: Uncertainty \& Conformal Prediction}}
\fancyhead[R]{\small 2026-07-08}
\fancyfoot[C]{\small\thepage}
\renewcommand{\headrulewidth}{0.8pt}

\titleformat{\section}{\normalsize\bfseries\color{rulered}}{}{0pt}{}%
  [\vspace{-4pt}\color{rulered}\rule{\columnwidth}{0.4pt}]
\titlespacing{\section}{0pt}{8pt}{4pt}

\hypersetup{colorlinks=true,urlcolor=rulered,linkcolor=black}

\begin{document}
\setlength{\parindent}{0pt}
\setlength{\parskip}{4pt}

{\small\color{keywordgray}\textbf{Keywords:}
  online conformal prediction \textbullet{}
  universal portfolio algorithms \textbullet{}
  regret bounds \textbullet{}
  adversarial streams}\par\vspace{4pt}

{\LARGE\bfseries\raggedright
  Online Conformal Prediction via Universal Portfolio Algorithms}\par\vspace{4pt}

{\large\itshape\raggedright
  A Parameter-Free Reduction Turns Coverage Tuning into Portfolio
  Selection}\par\vspace{6pt}

{\small\textbf{By} Tuo Liu, Edgar Dobriban \& Francesco Orabona
  (University of Pennsylvania, Boston University)}\par\vspace{2pt}

{\small\color{keywordgray}arXiv:2602.03168 \textbullet{}
  ICML 2026 Spotlight}
\par\vspace{2pt}\noindent{\color{rulered}\rule{\columnwidth}{0.6pt}}\vspace{6pt}

Online conformal prediction promises long-run coverage even on
adversarial data streams---but every prior method needed a hand-tuned
learning rate to actually deliver it. This paper shows that controlling
a general notion of \emph{linearized regret} implies coverage for
\emph{any} online algorithm, then reduces the resulting regret-minimization
problem to two-asset portfolio selection. The result, \textbf{UP-OCP}, is
parameter-free by construction, borrowing decades of tuning-free universal
portfolio algorithms to deliver calibrated coverage without a single
knob to turn.

\begin{mdframed}[backgroundcolor=boxgray,linecolor=rulered,linewidth=1pt,
    innerleftmargin=6pt,innerrightmargin=6pt,
    innertopmargin=6pt,innerbottommargin=6pt]
\textbf{\small AT A GLANCE}\par\vspace{4pt}
\begin{description}[leftmargin=0pt,labelsep=4pt,itemsep=2pt,parsep=0pt,
    font=\small\bfseries]
  \item[Problem:] Online conformal methods need a hand-tuned learning
    rate to guarantee coverage---exactly the kind of choice adversarial
    streams make impossible to validate in advance.
  \item[Why it matters:] Mistuned online conformal intervals silently
    lose their coverage guarantee on the non-stationary or adversarial
    streams they are meant to handle.
  \item[Method:] A regret-to-coverage reduction via the Fenchel
    conjugate, followed by a reduction of the resulting problem to
    two-asset portfolio selection.
  \item[Result:] UP-OCP matches or exceeds tuned ACI-style baselines'
    coverage with zero learning-rate parameters.
  \item[Evidence:] Comparisons against Adaptive Conformal Inference (ACI)
    variants on time-series and adversarial streams.
\end{description}
\end{mdframed}

\section{The Big Picture}

Online conformal prediction (OCP) methods maintain a calibrated
threshold as data arrive one at a time, aiming for long-run
$(1-\alpha)$ coverage even when the stream is non-stationary or
adversarial. The dominant family, descended from Adaptive Conformal
Inference (ACI), updates the threshold via an online-gradient-style
rule governed by a learning rate. That learning rate trades off how
quickly the method tracks distribution shift against how much it
overreacts to noise---and the correct value is stream-dependent and
unknowable in advance, precisely in the adversarial setting OCP exists
to handle without such foreknowledge.

This paper asks a structural question instead of a tuning question:
what property of an online algorithm's regret actually implies
coverage, for any algorithm at all?

\section{Why Existing Approaches Fall Short}

ACI-style methods require an algorithm-specific regret analysis for
every new update rule, and their coverage guarantees are only as good
as the learning rate chosen for that specific rule. Prior refinements
made the tuning more forgiving but never removed it. Because the
correctness of the guarantee depends on a hyperparameter that cannot be
cross-validated on an adversarial stream without assuming away the very
adversariality the method is supposed to survive, tuning is not merely
inconvenient---it is a genuine gap in the coverage guarantee's
practical reliability.

\section{Core Mechanism}

The authors show that controlling \emph{linearized regret} for the
$(1-\alpha)$-pinball loss implies a long-run coverage bound for any
online algorithm, via a black-box reduction that depends only on the
Fenchel conjugate of an upper bound on that regret. This decouples
``does it cover'' from ``which algorithm produced the threshold.''
Recognizing that the threshold-update problem, viewed through this
regret lens, is structurally identical to sequential betting between
two assets (covered vs.\ not covered), they import universal portfolio
algorithms---originally designed for sequential investing with no
prior knowledge of asset returns---directly into the conformal setting.

\section{Technical Formulation}

Let $s_t$ be the nonconformity score at time $t$ and $q_t$ the
calibrated threshold, with pinball loss
\[
\rho_{1-\alpha}(u) =
\begin{cases}
  (1-\alpha)\,u & u \ge 0 \\
  -\alpha\,u & u < 0
\end{cases},
\qquad u = q_t - s_t.
\]
Define regret
\[
R_T = \sum_{t=1}^{T} \rho_{1-\alpha}(q_t - s_t)
  - \min_{q} \sum_{t=1}^{T} \rho_{1-\alpha}(q - s_t).
\]
Empirical coverage error is controlled by $R_T$ through a bound derived
from the Fenchel conjugate of the regret's upper envelope. Casting
$q_t$'s update as a bet fraction between a ``covered'' and ``not
covered'' asset reduces the problem to constant-rebalanced portfolio
selection, for which universal portfolio algorithms guarantee
$O(\log T)$ regret with no learning-rate parameter---translating,
through the regret-to-coverage bound, into parameter-free coverage.

\section{Learning or Inference Procedure}

\begin{enumerate}[itemsep=2pt,parsep=0pt]
  \item Observe nonconformity score $s_t$ at each round $t$.
  \item Map the threshold-selection problem to a bet fraction over a
    two-asset (covered / not covered) portfolio.
  \item Update the bet using a universal portfolio algorithm---no
    learning rate is set or tuned.
  \item Convert the resulting bet fraction back into the calibrated
    threshold $q_{t+1}$ and output the prediction interval.
\end{enumerate}

\section{What the Guarantee Says}

For any sequence of scores, including adversarially chosen ones, the
empirical long-run miscoverage rate deviates from $\alpha$ by an amount
controlled by $R_T/T \to 0$, where $R_T$ is the universal portfolio
algorithm's $O(\log T)$ regret. The bound holds without assuming any
distributional structure on the stream and without any user-chosen
learning rate.

\section{Experimental Findings}

On time-series and deliberately adversarial data streams, UP-OCP
matches or exceeds the empirical coverage of ACI variants that required
per-stream learning-rate tuning, while UP-OCP itself uses none---removing
an entire axis of hyperparameter search from online conformal
deployment.

\section{Ablations and Interpretation}

\textbf{Against best-tuned ACI:} Even when ACI's learning rate is
tuned in hindsight on each specific stream, UP-OCP's parameter-free
coverage remains competitive, showing little practical cost to removing
the tuning step.

\textbf{Against mistuned ACI:} When ACI's learning rate is chosen
without stream-specific tuning, its coverage degrades substantially;
UP-OCP's guarantee is unaffected, since it never depended on a tuned
parameter in the first place.

\textbf{Beyond intervals:} The authors note the regret-to-coverage
reduction is not intrinsically tied to the two-asset case, suggesting
the portfolio-algorithm connection may extend to richer conformal
prediction settings.

\section*{Reference}

Tuo Liu, Edgar Dobriban, Francesco Orabona.
\textbf{Online Conformal Prediction via Universal Portfolio Algorithms}.
arXiv:2602.03168, February 2026.
\url{https://arxiv.org/abs/2602.03168}

\end{document}
